% NichidaiMasterExam_R7_2nd_2
\begin{tcolorbox} %R7_2nd_2
	$\fbox{2}$ $n$次実正方行列$A$が次の条件をみたすとき、$A$を正定値行列という。ただし$\tr{\bm{v}}$は$\bm{v}$の転置とする。
\begin{tcolorbox}
	零ベクトルでない任意の$n$次元実数ベクトル$\bm{v}$に対し、$\tr{\bm{v}} A\bm{v}> 0$
\end{tcolorbox}
	$H_n$を$i$行$j$列の要素が$\frac{1}{i+ j- 1}$である$n$次実正方行列とする。
\begin{enumerate}
	\item $H_3$の階数($\rank{}$)を求めよ。
	\item $H_2$は定数値行列であることを示せ。
	\item $n$次元ベクトル$\bm{v}= \tr{(v_1\ v_2\cdots\ v_n)}$に対し、以下を示せ。
\begin{equation*}
	\tr{\bm{v}} H_n\bm{v}= \sum_{i= 1}^n\sum_{j= 1}^n \frac{v_i v_j}{i+ j-1}
\end{equation*}
	\item $n$次元ベクトル$\bm{v}= \tr{(v_1\ v_2\cdots\ v_n)}$に対し、多項式$f_v(x)$を次のように定める
\begin{equation*}
	f_v(x)= v_n x^{n-1}+ v_{n- 1} x^{n- 2}+ \cdots+ v_2 x+ v_1
\end{equation*}
	以下を示せ。
\begin{equation*}
	\tr{\bm{v}} H_n\bm{v}= \int_0^1 {f_v(x)}^2 dx
\end{equation*}
	\item $H_n$が正定値行列であることを示せ。
\end{enumerate}
\end{tcolorbox}

\begin{souanswer} %%R7_2nd_2_ans
	$\fbox{2}$
\begin{enumerate}
	\item $H_3$は以下のように書き下せる。
\begin{equation*}
	H_3= 
\begin{pmatrix}
	1& 1/2& 1/3\\
	1/2& 1/3& 1/4\\
	1/3& 1/4& 1/5
\end{pmatrix}
\end{equation*}
	行列式を調べると、
\begin{equation*}
	\begin{vmatrix}H_3\end{vmatrix}= 1\cdot \left(\frac{1}{15}- \frac{1}{16}\right)- \frac{1}{2}\left(\frac{1}{10}- \frac{1}{12}\right)+ \frac{1}{3}\left(\frac{1}{8}- \frac{1}{9}\right)= \frac{1}{240}- \frac{1}{120}+ \frac{1}{216}= \frac{1}{21600}\ne 0
\end{equation*}
	$\begin{vmatrix}H_3\end{vmatrix}\ne 0$なのでフルランク。$\rank{H_3}= 3$

	\item $H_2= \begin{pmatrix}1& 1/2\\ 1/2& 1/3\end{pmatrix}$で任意の非零ベクトル$\bm{v}= \tr{\begin{pmatrix}v_1& v_2\end{pmatrix}}$に対して
\begin{align*}
	\tr{\bm{v}}H_2\bm{v}
	&= \begin{pmatrix}v_1\\ v_2\end{pmatrix}\begin{pmatrix}1& 1/2\\ 1/2& 1/3\end{pmatrix}\begin{pmatrix}v_1& v_2\end{pmatrix}\\
	&= v_1^2+ v_1v_2+ \frac{1}{3}v_2^2\\
	&= \left(v_1+ \frac{1}{2}v_2\right)^2+ \frac{1}{12}v_2^2
\end{align*}
	となる。$v_1, v_2$のうち少なくとも一方が0でなければ、この値は常に正となる。すなわち$H_2$は正定値行列。

	\item $\bm{v}= (v_1, \ldots, v_n)^\mathrm{T}$、および$H_n$の要素$(H_n)_{ij}= \frac{1}{i+ j- 1}$より
\begin{align*}
	\bm{v}^\mathrm{T}H_n\bm{v}
	&= \sum_{i= 1}^n v_i (H_n)_i\\
	&= \sum_{i= 1}^n v_i\left(\sum_{j= 1}^n \frac{1}{i+ j- 1}v_j\right)\\
	&= \sum_{i= 1}^n \sum_{j= 1}^n\frac{v_i v_j}{i+ j- 1}
\end{align*}

	\item $f_v(x)= \sum_{k= 1}^n v_kx^{k- 1}$とおくと
\begin{align*}
	\{f_v(x)\}^2
	&= \left(\sum_{i= 1}^n v_ix^{i- 1}\right)\left(\sum_{j= 1}^n v_jx^{j- 1}\right)\\
	&= \sum_{i= 1}^n\sum_{j= 1}^n v_iv_j x^{i+ j- 2}
\end{align*}
	これを0から1まで積分する。
\begin{equation*}
	\int_0^1\{f_v(x)\}^2 dx= \sum_{i= 1}^n \sum_{j= 1}^n v_iv_j \int_0^1 x^{i+ j-2} dx
\end{equation*}
	ここで$\int_0^1 x^{i+ j- 2} dx= \frac{1}{i+ j- 1}$なので、
\begin{equation*}
	\int_0^1\{f_v(x)\}^2 dx= \sum_{i= 1}^n \sum_{j= 1}^n\frac{v_iv_j}{i+ j- 1}
\end{equation*}

	\item 任意の$\bm{v}\ne 0$に対して、$\bm{v}^\mathrm{T}H\bm{v}= \int_0^1 \{f_v(x)\}^2 dx$
\begin{itemize}
	\item $\{f_v(x)\}^2\ge 0$であるからその積分値は0以上
	\item  $\int_0^1 \{f_v(x)\}^2 dx= 0$となるのは連続関数$f_v(x)$が区間$[0, 1]$で恒等的に0である時に限られる。
	\item $f_v(x)$は$n- 1$次多項式であり、$\bm{v}\ne 0$であれば、係数の少なくとも1つが0ではないため、$f_v(x)$は恒等的に0にはならない。
\end{itemize}
\end{enumerate}
\end{souanswer}